\documentclass[10pt,a4paper]{amsart}
\usepackage[margin=19mm]{geometry}
\usepackage{amsmath,amssymb,mathtools}
\usepackage[scr=rsfs,cal=euler]{mathalfa}
\usepackage{xfrac}
\usepackage[unicode,hidelinks,bookmarks=false]{hyperref}
\newcommand{\ie}{i.e.\ }
\newcommand{\cf}{cf.\ }
\newcommand{\resp}{respectively\ }
\newcommand{\Subsection}{\subsection}
\providecommand{\nobreakdash}{\nobreak\hbox{-}}
\setlength{\emergencystretch}{2em}
\multlinegap0pt
\usepackage{amssymb}
\usepackage{float}
\usepackage[all]{xy}
\usepackage{cancel}
\newtheorem{theorem}{Theorem}
\newtheorem{corollary}{Corollary}
\newcommand{\Ham}{\operatorname{Ham}}
\newcommand{\Hcal}{{\mathcal{H}}}
\title{On the intersections of projected Hamiltonian orbits in cotangent bundles}
\author{Lucas Dahinden, Jacobus de Pooter}
\date{Source: arXiv:2602.15693v1 (CC BY 4.0)}
\begin{document}
\maketitle

\noindent\emph{Source and licence.} On the intersections of projected Hamiltonian orbits in cotangent bundles. Version arXiv:2602.15693v1 (CC BY 4.0), distributed by arXiv under the Creative Commons Attribution 4.0 International licence, http://creativecommons.org/licenses/by/4.0/, as recorded in the arXiv licence metadata for this exact version and shown on the abstract page https://arxiv.org/abs/2602.15693v1. Full licence notice: \texttt{licenses/arxiv-2602.15693-CC-BY-4.0.txt}.\par\smallskip

\noindent\textit{Editorial scope.} Counted unit: the article's corollary cor:dimensioncount (source0.tex lines 1046--1048). The article's complete original proof of it is delivered with it (source0.tex lines 1049--1059, one \texttt{proof} environment of 11 lines). No mathematics is written, repaired or re-derived: the statement and the proof are the article's own lines, in the article's own notation, and no retained line is substituted or re-flowed.  Retained uncounted as the source-local context the proof needs: lines 992--998.  Nothing was omitted from any retained span, so the package carries no omission record.  No retained line contains a citation, so the wrapper carries no bibliography and no bibliography entry.  No retained line contains a figure environment, an image-inclusion command or drawing code, so no image is delivered and the package declares no asset dependency.  Editorial formatting only, in addition: an AMS \texttt{amsart} preamble that restates the article's own declarations and macros used by the retained lines, namely the environments \texttt{theorem}, \texttt{corollary} on separate counters, together with the standard packages those lines need.  The retained environments print in this document as Theorem 1, Corollary 1 in order of appearance, whereas the article prints the same items with the numbering of its own full document.  The complete native source of the article as distributed in the arXiv e-print for this exact version is bundled unchanged as \texttt{sources/194814/source0.tex}.
\begin{theorem}[Transversality for homopodal points]\label{thm:homopodaltransversality}
    Let $Q$ be a manifold of dimension $d\geq 2$. Let $k\geq 1$. For a $C^{k+2}$-residual set of Hamiltonians $H\in \Ham_s(Q)$, the following is true:
    \begin{itemize}
        \item The set of antipodal points $\Hcal^-_k(\Sigma)\subseteq \Sigma\times\Sigma$ is a submanifold of dimension $(3-k)(n-1)+1$.
        \item The set of non-diagonal isopodal points $\Hcal^+_k\setminus \Delta_\Sigma\subseteq \Sigma\times\Sigma$ is a submanifold of dimension $(3-k)(n-1)+1$. 
    \end{itemize}
\end{theorem}

\begin{corollary}\label{cor:dimensioncount}
    If $n=2$ and $k\geq 5$, or if $n\geq 3$ and $l\geq 4$, then there is a $C^{k+2}$-residual subset of $\Ham_s$ such that the set of $k$-th order homopodes is empty except for the diagonal $\Hcal_k\setminus\Delta_{\Sigma}=\emptyset$.
\end{corollary}

\begin{proof}
    We merely need to estimate the dimension $(3-k)(n-1)+1$. If the dimension is negative, then $\Hcal_k$ is empty except for the diagonal. By Theorem~\ref{thm:homopodaltransversality} transversality can be arranged for a $C^{k+2}-$residual subset of $\Ham_s$.
    
    If $n=2$, then 
    $$(3-k)(n-1)+1=4-k,$$
    which is negative for $k\geq 5$. 

    If $n\geq 3$ and $k\geq 4$, then 
    $$(3-k)(n-1)+1 \leq -1,$$
    which is negative. 
\end{proof}

\end{document}
